\paragraph{\trLagrange\trBrane。}\label{sub:lagrangian-branes}
设 $(M,\varphi_M)$ 是一个满足 $2c_1(M)=0$ 的\trSymplectic\trManifold，并配备一个%
\trCompatible{}的\trAlmostComplex{}结构 $I_M$。取定一个二次复\trVolume{}形式
$\eta_M^2$，其含义是（按假设为平凡的）双\trCanonical{}丛
$\scr{K}_M^2=\lambda_{\mathbb{C}}^{\mathrm{top}}(TM)^{\otimes-2}$ 的一个无处为零的%
\trSection。令 $\alpha_M\colon Gr(TM)\to S^1$ 为与之相伴的\trSquaredPhase{}映射，见
\eqref{eq:squared-phase-map}。一个\textbf{\trLagrange\trBrane}（Lagrangian
brane）是一个三元组
\begin{equation*}
L^{\#}=(L,\alpha^{\#},P^{\#})
\end{equation*}
由一个\trLagrange\trSubmanifold{}$L\subset M$ 与\trSubbundle{}$TL\subset TM|L$ 上的一个%
\trBraneStructure{}$(\alpha^{\#},P^{\#})$ 组成。具体地说，它由一个满足
$\exp(2\pi i\,\alpha^{\#}(x))=\alpha_M(TL_x)$ 的函数 $\alpha^{\#}\colon L\to\RR$
（即 $L$ 的\trGrading）连同 $TL$ 的一个\trPinStructure{}$P^{\#}$ 构成。这类结构的
基本理论就是\cref{sub:branes-revisited} 中一般讨论的一个特例。为了把它与%
\trSymplectic\trTopology{}中更常用的概念联系起来，应当指出：\trMaslovClass{}
$\mu_L\in H^1(L)$ 是\trGrading{}存在性的\trObstruction，它在边缘映射
$H^1(L)\to H^2(M,L)$ 下映到%
\trRelative\trFirstChernClass{}$2c_1(M,L)$。另一方面，$\mu_L$ 的 mod 2 约化是
$w_1(L)-w_1(\xi_M|L)$，其中 $\xi_M\subset\lambda_{\mathbb{C}}^{\mathrm{top}}(TM)$
是使 $\pm\eta_M$ 取实值的那个\trRealLineBundle。这与\cref{thm:grading-relative-orientation}
密切相关：后者说，一切\trLagrange\trBrane{}$L^{\#}$ 都带有相对于 $\xi_M|L$ 的%
\trCanonical\trOrientation；在 $\xi_M$ 平凡的特殊情形——这等于说 $\eta_M^2$ 有一个
整体定义的平方根 $\eta_M$——\trLagrange\trBrane{}就是寻常意义下\trOrientation{}的
（更确切地说，平方根的选取固定了诸\trOrientation）。我们把用\trRealLineBundle{}
$\beta\to L$ 改变\trPinStructure{}$P^{\#}$ 的运算记作 $L^{\#}\otimes\beta$，而把
$\sigma$ 重的\trSuspension{}操作记作 $S^\sigma L^{\#}$，后者在每一点处都由
\eqref{eq:shift-operation} 给出。
\begin{remark}\label{thm:brane-homotopy-remark}
\trLagrange\trBrane{}的概念本质上只依赖于 $\eta_M^2$ 的\trHomotopy{}类，这相当于
给出了实效为 $H^1(M)$ 的自由度。除此之外，该定义还有两种变体。其一，可以取定某个类
$d\in H^1(M;\ZZ/2)$，用一条\trRealLineBundle{}$\delta$ 来表示它，并把 $TL$ 上的%
\trPinStructure{}换成 $TL\oplus\delta|L$ 上的。其二，给定某个 $w\in H^2(M;\ZZ/2)$，
可以按\cref{thm:twisted-pin-structures}的意义考虑相对于 $w|L$ 的\trTwisted
\trPinStructure（这两种想法的组合也是可能的）。对此我们不再深究，有兴趣的读者请
转阅 [60]。不过，我们确实想概述一下：如何通过把\cref{sub:abstract-brane-structures}
中的构造全局化，设计出一个统一的框架来囊括所有这些可能性。
给定\trCohomology{}类 $\mu_M\in H^1(Gr(TM))$ 与 $w_M\in H^2(Gr(TM);\ZZ/2)$，二者逐%
\trFiber{}地限制为\trMaslovClass{}与第二\trStiefelWhitney{}类；于是，如
\eqref{eq:maslov-w2-classifying-map} 那样，可以构造一个\trFibration{}
$Gr(TM)^{\#}\to Gr(TM)$，其\trFiber{}为 $\ZZ\times\RP^\infty$。所谓\textbf{抽象%
\trLagrange\trBrane}（abstract Lagrangian brane），指的就是\trLagrange\trSubmanifold{}
$L\subset M$ 连同\trCanonical\trSection{}$L\to Gr(TM)|L$ 到 $Gr(TM)^{\#}$ 上的一个%
\trLift。为了与我们先前的术语挂钩，注意\trSquaredPhase{}映射在二次\trVolume{}形式
$\eta_M^2$ 的\trHomotopy{}类与\trCohomology{}类 $\mu_M$ 之间建立起\trBijective{}对应。
类似地，最一般的类 $w_M$ 具有形式 $w_2(F)+w_1(F)\cdot\pi^*d+\pi^*w$，其中
$F\to Gr(TM)$ 是\trUniversal{}丛，$\pi\colon Gr(TM)\to M$ 是\trProjection，而
$(d,w)$ 如前。
\end{remark}
